\section{纽结：Wilson 线与穿刺点}

\subsection{纽结从哪里进入}

第三、四讲处理的是一个没有纽结的三维流形：沿 $\Sigma$ 切开，得到一个干净的模空间。
现在要把一条 Wilson 线 $C$ 放进去。做法仍然是沿 $\Sigma$ 切开，但要选 $\Sigma$ 与 $C$ 横截相交；
于是 $\Sigma$ 上出现有限个交点 $P_1,\dots,P_n$，每个点带一个表示 $R_s$。

沿着 $C$ 的 Wilson 线 (2.5) 是“沿一条曲线的和乐”。在切开之后，这条曲线的两个端点
落在 $\Sigma$ 上，而中间部分在 $\Sigma$ 的一侧。如果我们只关心 $\Sigma$ 这一侧的量子态，
那么 Wilson 线的作用就是：在 $P_s$ 处插入一个“带表示 $R_s$ 的算符”。

\subsection{穿刺点上的高斯定律}

第三讲里高斯定律是 $F=0$。现在 $F$ 在穿刺点处不能为零：Wilson 线在那里是奇异的。
对 $A_0$ 变分，得到的约束变成
\begin{equation}
\frac{k}{8\pi}\,\varepsilon^{ij}F^a_{ij}(x)=\sum_s\delta^2(x-P_s)\,T^a_{(s)}, \tag{5.1}
\end{equation}
其中 $T^a_{(s)}$ 是表示 $R_s$ 的生成元。也就是说：曲率在 $P_s$ 处有一个“源”，
源的强度由表示决定。这正是电磁学里“点电荷给出电通量”的规范场版本。

\subsection{为什么不能经典地解}

(5.1) 看起来像一个方程，但它不是。原因有两条。

\textbf{第一，右端是算符。} $T^a_{(s)}$ 作用在表示空间 $R_s$ 上；不同的 $T^a$ 之间不对易。
因此 (5.1) 的右端不是一个数，而是一个算符值的分布。要“解”这个方程，意味着找到一个
数值联络 $A$ 使得两边相等，但右边的算符可以在 $R_s$ 上取不同的本征值，除非 $R_s$ 是一维的，
否则不存在一个同时满足所有分量的经典解。

\textbf{第二，$A$ 在 $P_s$ 处奇异。} 即使忽略算符问题，来源项也意味着 $A$ 在 $P_s$ 附近
不再是光滑的；平坦性只在 $\Sigma\setminus\{P_s\}$ 上成立。相空间因此是“带奇点的平坦联络模空间”，
而不是第三讲那个光滑的 $\cM$。

这两条合起来说明：穿刺点上的约束不是一个可以代入求解的方程，
而是一个必须在量子层面处理的“源条件”。

\subsection{把表示也量子化}

Witten 的处理办法是换一个视角：与其去解 (5.1)，不如把表示 $R$ 本身看成一个量子系统。
具体地，$R$ 是紧致群 $G$ 的不可约表示，而每一个不可约表示都可以实现为某个辛流形的几何量子化：
取余伴随轨道，或等价地取旗流形 $G/T$（$T$ 是极大环面）。$G/T$ 上有自然的 Kähler 结构，
对支配整权 $\lambda$，取辛形式 $\omega_\lambda$，其 de Rham 类就是 $\lambda$。
Borel--Weil--Bott 定理说：
\begin{equation}
H^0(G/T,\,\cL_\lambda)=R_\lambda, \tag{5.2}
\end{equation}
即 $G/T$ 上全纯线丛 $\cL_\lambda$ 的全纯截面空间，正是最高权为 $\lambda$ 的不可约表示。

注意这里 $\lambda$ 没有加 $\rho$：这是指标定理的直接结论。$\rho$ 的位移属于半形式修正
$K_{G/T}^{1/2}=\cL_{-\rho}$，不在 (5.2) 里。如果读者在自己的计算里看到 $\lambda+\rho$，
那多半是在用 Weyl 维数公式的写法，或者已经把半形式算进去了。

BWB 的证明思路值得知道，因为它解释了“为什么恰好是全纯截面”。$G/T$ 上的线丛 $\cL_\lambda$
的截面可以看成 $G$ 上的函数 $f$，满足 $f(gb)=\chi_\lambda(b)^{-1}f(g)$（$b\in B$，Borel 子群）。
按 Peter--Weyl 定理，这样的函数由最高权不超过 $\lambda$ 的表示张成；再要求全纯（即被
$\bar\partial$ 杀掉），就只剩下最高权恰为 $\lambda$ 的那个表示。于是 (5.2) 成立。
这条路线的关键是：不可约表示可以用“旗流形上的全纯截面”这个几何对象来实现，
而几何对象自然带辛结构，因而可以作为一个量子系统的相空间。

\subsection{抛物结构与扩大后的相空间}

把 (5.2) 代回，我们得到一个新的相空间：在 $\Sigma\setminus\{P_s\}$ 上取平坦 $G$-联络，
并要求在每个 $P_s$ 处，联络的结构群约化到 $T$（这样“源”就由 $\lambda_s$ 的整性给出）。
这个相空间记作 $\cM_{\{P_s,R_s\}}$。用丛的语言，这是在 $\Sigma$ 上取一个抛物 $G$-丛：
在 $P_s$ 处给定一个 Borel 子群（或一个滤过），平坦联络在这个滤过上取特定的形式。
抛物丛的模空间同样是 Kähler 的，只是多了边界（对应抛物权重的选择）。

在这个空间上，约束 (5.1) 变成了一个普通的经典约束（因为右端现在由 $G/T$ 的量子化提供，
不再是一个未定的算符），而辛形式也相应地扩大。把扩大后的相空间量子化，
得到的 Hilbert 空间记作 $\cH_{\Sigma;\{P_s,R_s\}}$。Witten 的论断是：
这个空间恰好是带有穿刺点 $P_s$、表示 $R_s$ 的共形块空间。
用第六讲的语言，它就是 WZW 模型的共形块。

\subsection{可积表示}

并非所有表示都给出非零的截面。只有“可积”的表示才有非零贡献；对 $SU(2)$，条件是最高的
自旋满足
\[
2j\le k.
\]
这个条件不是 Borel--Weil--Bott 的推论（BWB 对任意支配整权都成立），它来自仿射李代数：
level $k$ 的表示必须满足 $\lambda$ 在仿射 Weyl 群的基本区域里，而 $2j\le k$ 是它的
$SU(2)$ 特例。不可积的表示在量子化后给出零空间——用物理的语言说，这些表示“退耦”了。

\subsection{这是纲领，不是定理}

需要强调：上面这一段是 Witten 给出的\textbf{纲领}。原文用 "should" 描述它，
并明确说 "deferring a fuller treatment"——完整的构造留待后来。数学上把它严格化，
需要抛物丛、相交上同调、以及共形块理论；这些工作在 1990 年代之后才陆续完成。
本讲义把它作为理解的框架，而不是作为已证明的定理。

\subsection*{练习}
\begin{enumerate}[leftmargin=*]
  \item 说明 (5.1) 右端为什么是算符，并给出一个 $R$ 为一维表示时它可以经典求解的理由。
  \item 对 $G=SU(2)$、$R$ 为自旋 $j$ 的表示，写出 $\lambda$ 与 $j$ 的关系，并验证 BWB 给出的维数 $2j+1$。
  \item 为什么可积性条件是 $2j\le k$ 而不是 $2j\le k+2$？（提示：考虑半形式修正后哪个 $k$ 出现在条件里。）
\end{enumerate}

\begin{quote}
\small
\textbf{本讲要点。} Wilson 线在穿刺点处给出算符值的高斯定律 (5.1)；
Witten 用 $G/T$ 的几何量子化（BWB）把表示变成相空间的一部分，从而把约束经典化；
量子化后得到带穿刺点的共形块。可积性 $2j\le k$ 来自仿射代数。
\end{quote}
